\section{Prior work and the boundary of the present contribution}
\label{supp:prior-work}

The present contribution is a claim-specific computational reassessment,
not discovery of shrimp morphology, symbolic coding, or homoclinic hub
organization in general. Table~\ref{tab:prior-work} identifies the closest
prior objects and what our computations add or leave open. This is a targeted
comparison, not an exhaustive literature survey. The repository reference
ledger distinguishes full-text examination from abstract-level comparison;
unread proof and implementation details are not claimed as checked.

\small
\renewcommand{\arraystretch}{1.15}
\begin{longtable}{@{}P{0.24\textwidth}P{0.32\textwidth}P{0.35\textwidth}@{}}
\caption{Source-level comparison. Similar pictures, folds, or words need
not describe the same invariant object or parameter convention.}
\label{tab:prior-work}\\
\toprule
Prior work & Object or result & Consequence for this reassessment \\
\midrule
\endfirsthead
\toprule
Prior work & Object or result & Consequence for this reassessment \\
\midrule
\endhead
\bottomrule
\endfoot
Bonatto and Gallas (2008) \citep{bonatto2008periodicity} &
Numerical periodicity hubs and nested spirals in a resistive-circuit model. &
Hub morphology is prior work. This citation is not presented as an
experimental observation or as a new result of our atlas. \\[0.5em]
Letellier et al. (1995) \citep{letellier1995unstable} &
R{\"o}ssler unstable periodic orbits, symbolic descriptions, pruning, and templates. &
Our proposed partition must be compared with that established framework;
recovering a local alphabet alone is not a new template construction. \\[0.5em]
Barrio et al. (2009) \citep{barrio2009qualitative} &
Limit-cycle bifurcations and global R{\"o}ssler parameter-space organization. &
Orbit continuation is a methodological predecessor. Our corrected local
flip locus is not a claim to first mapping R{\"o}ssler bifurcations. \\[0.5em]
Vitolo et al. (2011) \citep{vitolo2011global} &
Global hub structure related to folded sheaves of Shilnikov homoclinic
bifurcations in a R{\"o}ssler variant. &
A folded homoclinic structure is not the folded period-6 flip locus.
Their $\dot z=(b+z)x-cz$ differs from Eq.~\eqref{eq:rossler}; equal symbols
do not justify comparing numerical coordinates directly. \\[0.5em]
Barrio et al. (2011, 2012, 2013)
\citep{barrio2011global,barrio2012topological,barrio2013spirals} &
Shilnikov hub organization, a topology-change construction, and double
superstability in a scalar description. &
The 2011 foundation remains explicit. Our finite classifier bracket is
not yet a reconstruction of an invariant topology-transition curve. \\[0.5em]
Jones (2012) \citep{jones2012topological} &
Finite symbolic ordering and reinjection-based period-increment links. &
The 2012 connection is treated as an independent, near-simultaneous
co-discovery with Barrio--Blesa--Serrano. The specific words and arrows,
not generic doubling, are our still-open reproduction target. \\[0.5em]
Malykh et al. (2020) \citep{malykh2020homoclinic} &
Homoclinic chaos, tailored symbolic analysis, and interaction of saddle-foci. &
The second equilibrium is a possible global contributor, not excluded by
our small-equilibrium candidate. Parameter conventions require translation. \\[0.5em]
Gierzkiewicz and Zgliczy\'nski (2021); Galias (2025)
\citep{gierzkiewicz2021periodic,galias2025symbolic} &
Computer-assisted periodic-orbit results; symbolic search, continuation,
and interval proofs of many periodic windows near the classical case. &
High precision and small residuals do not supply these existence guarantees.
Our distinctive aim is testing a specified hub mechanism, not maximizing
the number of detected periods or windows. \\[0.5em]
Xing and Luo (2023) \citep{xing2023period1} &
Semi-analytical period-1 motions toward approximate twin spiral homoclinic
orbits. &
An alternative orbit-to-homoclinic treatment; not a validation of our
printed-coordinate endpoint or of its organizing role. \\[0.5em]
Igra (2025) \citep{igra2025knots} &
Analytical knot/chaos results for an idealized model under trefoil and
heteroclinic assumptions. &
Those assumptions have not been established at our hub. The result cannot
be transferred to our numerical candidate by a visual resemblance. \\[0.5em]
do Carmo et al. (2025) \citep{docarmo2025measure} &
Local-maximum return-map density fits near hubs in five three-dimensional systems. &
Cross-system return-measure comparison is prior work, distinct from
replicating a particular symbolic reinjection connection. \\
\end{longtable}
\normalsize

The strongest prospective mechanism contribution would join the currently
separate observations: persistence of recovered unstable orbits, the
finite-sample lobe association, a defined invariant manifold/pruning event,
and a predicted symbolic insertion. No claim of a new completed mechanism
is made before that chain is tested. A visually folded curve or successful
symbolic encoding does not remove the need to identify the mathematical
object and its domain.
